\documentclass[11pt,a4paper]{article}

\usepackage[utf8]{inputenc}
\usepackage[T1]{fontenc}
\usepackage[french]{babel}
\usepackage[margin=2.3cm]{geometry}
\usepackage{amsmath,amssymb}
\usepackage{booktabs}
\usepackage{array}
\usepackage{xcolor}
\usepackage{tikz}
\usetikzlibrary{arrows.meta,positioning,automata,calc}
\usepackage[hidelinks]{hyperref}

\definecolor{nettoyer}{RGB}{31,94,168}
\definecolor{charger}{RGB}{196,72,40}

\newcommand{\N}{\textsc{Nettoyer}}
\newcommand{\C}{\textsc{Charger}}
\newcommand{\pd}{p_{\text{dust}}}

\title{Case study 2 -- Le robot nettoyeur\\[2pt]
\large Modélisation et résolution par un processus de décision markovien}
\author{Lagier Hadrien, Serratos Paul, Tsan Alex}
\date{M2 -- Décision}

\begin{document}
\maketitle

\noindent Le modèle est implémenté en Python avec la bibliothèque \texttt{pymdptoolbox} de l'INRAE.
Le programme fourni (\texttt{iadec26\_2-MDPs\_CS-robot-nett-fr.py}) construit les matrices décrites
ci-dessous, vérifie qu'elles sont stochastiques et calcule les politiques optimales présentées dans ce rapport.

%==========================================================================
\section{Modélisation (question A)}
%==========================================================================

\subsection{États}

Un état combine la position du robot et la propreté du salon. Le robot lui-même n'est ni propre ni sale :
la seconde lettre décrit toujours l'état \emph{du salon}.

\begin{center}
\begin{tabular}{>{\bfseries}l l l}
\toprule
État & Position du robot & Salon \\
\midrule
LD & salon (\emph{Living room}) & sale (\emph{Dirty}) \\
LC & salon & propre (\emph{Clean}) \\
CD & station de charge (\emph{Charger}) & sale \\
CC & station de charge & propre \\
O  & hors tension (\emph{Off}) & -- \\
\bottomrule
\end{tabular}
\end{center}

\subsection{Actions}

Deux actions sont disponibles dans tous les états : \N{} et \C.
\begin{itemize}
  \item \N{} : le robot se rend (ou reste) dans le salon et le nettoie. Il arrive donc toujours dans un état L$\ast$.
  \item \C{} : le robot se rend (ou reste) à la station de charge. Il arrive donc toujours dans un état C$\ast$.
  \item En O, les deux actions ont le même effet : le robot reste hors tension.
\end{itemize}

\subsection{Hypothèses d'interprétation}

L'énoncé laisse quelques points ouverts. Nous les tranchons ainsi :
\begin{enumerate}
  \item \textbf{La poussière concerne le salon.} « Lorsque le robot charge, [le salon] peut passer de propre
        à sale avec une probabilité $\pd$ ». C'est cohérent avec la suite de la phrase (« si le salon est déjà sale,
        il reste sale »).
  \item \textbf{\N{} depuis la station} signifie quitter la station pour aller nettoyer le salon. Le robot part
        batterie pleine, donc ce nettoyage ne peut pas provoquer de panne. Depuis le salon, la batterie a déjà servi,
        d'où le risque de panne $e$.
  \item \textbf{O est absorbant} : une fois hors tension, le robot n'en sort plus.
  \item Le salon ne se salit que pendant la charge (seule source de poussière mentionnée par l'énoncé).
\end{enumerate}

\subsection{Paramètres}

\begin{center}
\begin{tabular}{>{$}c<{$} l}
\toprule
e & probabilité de tomber en panne en nettoyant depuis le salon \\
\varepsilon & probabilité d'échec du nettoyage d'un sol sale \\
\pd & probabilité que le salon se salisse pendant une charge \\
\gamma & facteur d'actualisation (horizon infini), $\gamma = 0{,}8$ \\
\bottomrule
\end{tabular}
\end{center}

\subsection{Graphe du MDP}

La figure~\ref{fig:mdp} représente toutes les transitions possibles. Les flèches pleines bleues
correspondent à l'action \N, les flèches pointillées rouges à l'action \C.

\begin{figure}[ht]
\centering
\begin{tikzpicture}[
    >={Stealth[length=2.6mm]},
    state/.style={circle, draw, thick, minimum size=13mm, font=\bfseries},
    n/.style={->, thick, nettoyer},
    c/.style={->, thick, charger, dashed},
    lbl/.style={font=\footnotesize, fill=white, inner sep=1.2pt},
    every loop/.style={looseness=6}
  ]
  \node[state] (LD) at (0,3.6) {LD};
  \node[state] (LC) at (7,3.6) {LC};
  \node[state] (CD) at (0,0)   {CD};
  \node[state] (CC) at (7,0)   {CC};
  \node[state, fill=black!8] (O) at (3.5,7) {O};

  % --- Nettoyer ---
  \draw[n] (LD) to node[lbl, pos=0.5] {$(1-e)(1-\varepsilon)$} (LC);
  \draw[n] (LD) to[out=150, in=210, loop] node[lbl, left] {$(1-e)\varepsilon$} (LD);
  \draw[n] (LD) to node[lbl, pos=0.45] {$e$} (O);
  \draw[n] (LC) to[out=30, in=-30, loop] node[lbl, right] {$1-e$} (LC);
  \draw[n] (LC) to node[lbl, pos=0.45] {$e$} (O);
  \draw[n] (CD) to[bend left=22] node[lbl, pos=0.5] {$\varepsilon$} (LD);
  \draw[n] (CD) to[bend left=8] node[lbl, pos=0.3] {$1-\varepsilon$} (LC);
  \draw[n] (CC) to[bend right=22] node[lbl, pos=0.5] {$1$} (LC);

  % --- Charger ---
  \draw[c] (LD) to[bend left=22] node[lbl, pos=0.5] {$1$} (CD);
  \draw[c] (CD) to[out=210, in=150, loop] node[lbl, left] {$1$} (CD);
  \draw[c] (LC) to[bend right=22] node[lbl, pos=0.5] {$1-\pd$} (CC);
  \draw[c] (LC) to[bend left=8] node[lbl, pos=0.3] {$\pd$} (CD);
  \draw[c] (CC) to[out=-30, in=30, loop] node[lbl, right] {$1-\pd$} (CC);
  \draw[c] (CC) to node[lbl, pos=0.5] {$\pd$} (CD);

  % --- O absorbant ---
  \draw[n] (O) to[out=60, in=120, loop] node[lbl, above] {$1$ (les deux actions)} (O);

  % --- Légende ---
  \begin{scope}[shift={(9.6,6.6)}]
    \draw[n] (0,0) -- (1,0) node[right, black, font=\small] {\N};
    \draw[c] (0,-0.6) -- (1,-0.6) node[right, black, font=\small] {\C};
  \end{scope}
\end{tikzpicture}
\caption{Graphe du MDP du robot nettoyeur (5 états, 2 actions).}
\label{fig:mdp}
\end{figure}

\subsection{Matrices de transition}

Les états sont ordonnés (LD, LC, CD, CC, O). L'élément $(s, s')$ de $P^a$ vaut $P(s' \mid s, a)$ ;
chaque ligne somme à 1 (vérifié dans le programme par \texttt{verify\_transition\_add\_to1} et
\texttt{mdptoolbox.util.check}).

\[
P^{\N} =
\begin{pmatrix}
(1-e)\varepsilon & (1-e)(1-\varepsilon) & 0 & 0 & e \\
0 & 1-e & 0 & 0 & e \\
\varepsilon & 1-\varepsilon & 0 & 0 & 0 \\
0 & 1 & 0 & 0 & 0 \\
0 & 0 & 0 & 0 & 1
\end{pmatrix}
\qquad
P^{\C} =
\begin{pmatrix}
0 & 0 & 1 & 0 & 0 \\
0 & 0 & \pd & 1-\pd & 0 \\
0 & 0 & 1 & 0 & 0 \\
0 & 0 & \pd & 1-\pd & 0 \\
0 & 0 & 0 & 0 & 1
\end{pmatrix}
\]

Dans \texttt{mdptoolbox}, ces deux matrices forment un tableau $P$ de dimension
$(|A|, |S|, |S|) = (2, 5, 5)$.

\subsection{Matrice de récompense}

L'énoncé attache la récompense à l'\emph{état} : $r_c$ pour un salon propre, $r_d$ pour un salon sale,
$\mathit{costcrash}$ hors tension. La récompense ne dépend donc pas de l'action ; dans \texttt{mdptoolbox},
$R$ est de dimension $(|S|, |A|)$ avec deux colonnes identiques :
\[
R =
\begin{pmatrix}
r_d & r_d \\
r_c & r_c \\
r_d & r_d \\
r_c & r_c \\
\mathit{costcrash} & \mathit{costcrash}
\end{pmatrix}
\quad \text{(lignes LD, LC, CD, CC, O ; colonnes \N, \C)}.
\]

Nous avons volontairement écarté une récompense qui noterait les \emph{décisions}
(par exemple $+10$ pour « nettoyer un salon sale »). Une telle récompense incite le robot à laisser le salon se salir pour
pouvoir le nettoyer à nouveau : en simulation, avec $e = 0$ et $\pd = 0{,}9$, le salon n'était plus propre que
49\,\% du temps, contre 100\,\% avec une récompense par état. La récompense décrit \emph{l'objectif} ;
c'est à l'optimisation de trouver \emph{les actions} qui l'atteignent.

%==========================================================================
\section{Valeurs des récompenses (question B)}
%==========================================================================

\begin{center}
\begin{tabular}{>{$}c<{$} r l}
\toprule
r_c & $10$ & un salon propre est l'objectif du robot \\
r_d & $0$ & un salon sale est une situation neutre, à corriger \\
\mathit{costcrash} & $-20$ & une panne est pénalisée à chaque pas, pour toujours \\
\bottomrule
\end{tabular}
\end{center}

Ces valeurs respectent les contraintes de l'énoncé, $\mathit{costcrash} \le r_d < r_c$ (vérifiées dans le programme).

Comme O est absorbant, la pénalité de panne se cumule indéfiniment :
\[
V(\text{O}) = \sum_{t \ge 0} \gamma^t \cdot \mathit{costcrash} = \frac{-20}{1 - 0{,}8} = -100.
\]
Une panne coûte donc l'équivalent de dix pas passés dans un salon propre. C'est cet écart qui
rend le robot prudent : nous n'avons pas choisi $\mathit{costcrash}$ beaucoup plus négatif, sinon le robot
refuserait de prendre le moindre risque quelles que soient les valeurs de $e$.

%==========================================================================
\section{Résolution et résultats (question C)}
%==========================================================================

\subsection{Méthode}

Pour chaque scénario, nous construisons $P$ et lançons \texttt{mdptoolbox.mdp.ValueIteration(P, R, 0.8)}
(convergence en 43 itérations). Nous avons vérifié avec \texttt{PolicyIteration} que les deux algorithmes
donnent la même politique dans tous les cas. Les paramètres communs sont $\varepsilon = 0{,}2$ et $\gamma = 0{,}8$.

La valeur d'une action dans un état est
\[
Q(s,a) = R(s,a) + \gamma \sum_{s'} P(s' \mid s,a)\, V(s'),
\]
et la politique optimale choisit dans chaque état l'action de plus grande valeur $Q$.

\subsection{Politiques optimales}

\begin{center}
\begin{tabular}{l c c c c c c}
\toprule
Scénario & $e$ & $\pd$ & LD & LC & CD & CC \\
\midrule
1. Panne probable & $0{,}5$ & $0{,}3$ & \C & \C & \N & \N \\
2. Panne peu probable & $0{,}05$ & $0{,}3$ & \N & \C & \N & \N \\
3. Très bonne batterie & $0{,}001$ & $0{,}3$ & \N & \N & \N & \N \\
4. Salon se salit vite & $0{,}1$ & $0{,}9$ & \C & \C & \N & \N \\
\bottomrule
\end{tabular}
\end{center}

En O, les deux actions sont équivalentes ; la politique y est sans objet.

\begin{center}
\begin{tabular}{l r r r r r}
\toprule
Scénario & $V(\text{LD})$ & $V(\text{LC})$ & $V(\text{CD})$ & $V(\text{CC})$ & $V(\text{O})$ \\
\midrule
1 & 24,4 & 41,5 & 30,5 & 43,2 & $-100$ \\
2 & 25,1 & 41,6 & 30,6 & 43,3 & $-100$ \\
3 & 37,5 & 49,4 & 37,6 & 49,5 & $-100$ \\
4 & 15,6 & 26,5 & 19,4 & 31,2 & $-100$ \\
\bottomrule
\end{tabular}
\end{center}

\subsection{Interprétation}

\paragraph{Le chargeur comme « point de départ sûr ».}
Dans tous les scénarios, depuis la station (CD, CC) le robot choisit \N. Ce nettoyage est sans risque de panne :
c'est la meilleure action possible. Le robot dispose donc toujours d'un chemin sûr vers un salon propre,
\[
\text{LD} \xrightarrow{\C} \text{CD} \xrightarrow{\N} \text{LC},
\qquad
\text{LC} \xrightarrow{\C} \text{CC} \xrightarrow{\N} \text{LC},
\]
qui ne coûte qu'un pas de temps supplémentaire. Toute la question est de savoir si ce détour vaut le risque
de panne évité.

\paragraph{Scénario 1 (panne probable).}
Avec $e = 0{,}5$, nettoyer depuis le salon revient à risquer une chance sur deux de perdre $V(\text{O}) = -100$.
Le robot ne nettoie donc jamais depuis le salon : il repasse systématiquement par la station
(cycle « recharger, puis repartir nettoyer »).

\paragraph{Scénario 2 (panne peu probable).}
Avec $e = 0{,}05$, la politique se scinde :
\begin{itemize}
  \item en LD, le robot nettoie directement : le gain d'un salon propre immédiatement l'emporte sur 5\,\% de risque ;
  \item en LC, le robot va charger : le salon est déjà propre, il n'y a rien à gagner à prendre un risque.
\end{itemize}
Le robot n'accepte le risque de panne que lorsqu'il y a un bénéfice immédiat.

\paragraph{Scénario 3 (très bonne batterie).}
Avec $e = 0{,}001$, le risque est négligeable : le robot nettoie partout et reste dans le salon, qui ne se
salit jamais pendant qu'il y est. Le détour par la station ne présente plus d'intérêt, et il exposerait même le
salon à la poussière.

\paragraph{Scénario 4 (salon qui se salit vite).}
Avec $e = 0{,}1$, la politique est la même qu'au scénario 1 : le risque de panne domine. La poussière fait
cependant baisser toutes les valeurs (un salon propre se resalit dès que le robot va charger).
L'effet de $\pd$ sur la politique n'apparaît que lorsque $e$ est petit, comme le montre l'étude ci-dessous.

\subsection{Étude de sensibilité}

Pour mieux situer les changements de politique, nous avons fait varier $e$ pour deux valeurs de $\pd$
(politique en LD / LC ; en CD et CC, le robot nettoie toujours) :

\begin{center}
\begin{tabular}{c c c}
\toprule
$e$ & $\pd = 0{,}3$ & $\pd = 0{,}9$ \\
\midrule
$0{,}2$ et $0{,}1$ & \C{} / \C & \C{} / \C \\
$0{,}05$ & \N{} / \C & \C{} / \C \\
$0{,}03$ à $0{,}015$ & \N{} / \C & \N{} / \N \\
$\le 0{,}01$ & \N{} / \N & \N{} / \N \\
\bottomrule
\end{tabular}
\end{center}

Deux effets apparaissent :
\begin{itemize}
  \item \textbf{Un salon qui se salit vite pousse le robot à rester dans le salon} quand la batterie est
        fiable : pour $e$ entre $0{,}015$ et $0{,}03$, LC passe de \C{} à \N{} quand $\pd$ augmente.
        Aller charger, c'est laisser la poussière s'installer.
  \item \textbf{Mais il rend le robot plus prudent quand le risque est modéré} : pour $e = 0{,}05$, LD passe de
        \N{} à \C. Un salon propre a moins de valeur s'il se resalit vite, donc le gain à le nettoyer
        immédiatement pèse moins face au risque de panne.
\end{itemize}

%==========================================================================
\section{Bonus : plusieurs niveaux de batterie}
%==========================================================================

\subsection{Modèle}

Nous ajoutons un niveau de batterie $b \in \{\text{H}, \text{M}, \text{L}\}$ (haut, moyen, bas) à chaque état.
On obtient $2 \times 2 \times 3 = 12$ états LDH, LDM, \dots, CCL, plus l'état O. Un seul état hors tension suffit,
puisque « hors tension » signifie batterie vide : il y a donc \textbf{13 états}.

Le risque de panne ne dépend plus de la position du robot mais de son niveau de batterie, ce qui est plus réaliste.
Le paramètre $e$ est remplacé par $e_b$ :
\[
e_{\text{H}} = 0{,}01, \qquad e_{\text{M}} = 0{,}1, \qquad e_{\text{L}} = 0{,}25.
\]

\begin{center}
\begin{tabular}{l l l l}
\toprule
Action & Position & Batterie & Salon \\
\midrule
\N & $\to$ salon & panne avec prob.\ $e_b$, sinon perd un niveau & sale $\to$ propre avec prob.\ $1-\varepsilon$ \\
\C & $\to$ station & gagne un niveau & propre $\to$ sale avec prob.\ $\pd$ \\
\bottomrule
\end{tabular}
\end{center}

Le niveau bas reste bas en nettoyant, et le niveau haut reste haut en chargeant.
Par exemple, depuis un état de batterie $b$ avec un salon sale :
\[
P(\text{LC}b^- \mid \text{xD}b, \N) = (1-e_b)(1-\varepsilon), \quad
P(\text{LD}b^- \mid \text{xD}b, \N) = (1-e_b)\,\varepsilon, \quad
P(\text{O} \mid \text{xD}b, \N) = e_b,
\]
où $b^-$ est le niveau inférieur à $b$ et x la position (L ou C). Les récompenses suivent la même règle que
le modèle de base : $r_c$ pour les six états « propre », $r_d$ pour les six états « sale », $\mathit{costcrash}$ pour O.

La figure~\ref{fig:batt} montre l'évolution du niveau de batterie selon l'action.

\begin{figure}[ht]
\centering
\begin{tikzpicture}[
    >={Stealth[length=2.6mm]},
    state/.style={circle, draw, thick, minimum size=11mm, font=\bfseries},
    n/.style={->, thick, nettoyer},
    c/.style={->, thick, charger, dashed},
    lbl/.style={font=\footnotesize, fill=white, inner sep=1.2pt},
    every loop/.style={looseness=6}
  ]
  \node[state] (H) at (0,0) {H};
  \node[state] (M) at (3.5,0) {M};
  \node[state] (L) at (7,0) {L};
  \node[state, fill=black!8] (O) at (3.5,-3.2) {O};

  \draw[n] (H) to[bend left=25] node[lbl] {$1-e_{\text{H}}$} (M);
  \draw[n] (M) to[bend left=25] node[lbl] {$1-e_{\text{M}}$} (L);
  \draw[n] (L) to[out=30, in=-30, loop] node[lbl, right] {$1-e_{\text{L}}$} (L);
  \draw[n] (H) to node[lbl] {$e_{\text{H}}$} (O);
  \draw[n] (M) to node[lbl] {$e_{\text{M}}$} (O);
  \draw[n] (L) to node[lbl] {$e_{\text{L}}$} (O);

  \draw[c] (L) to[bend left=25] node[lbl] {$1$} (M);
  \draw[c] (M) to[bend left=25] node[lbl] {$1$} (H);
  \draw[c] (H) to[out=150, in=210, loop] node[lbl, left] {$1$} (H);
\end{tikzpicture}
\caption{Évolution du niveau de batterie (la position et la propreté du salon évoluent comme dans le modèle de base).}
\label{fig:batt}
\end{figure}

\subsection{Résultats}

Avec $\varepsilon = 0{,}2$ et $\gamma = 0{,}8$, nous avons testé trois scénarios :
\begin{itemize}
  \item batterie standard : $e_b = (0{,}01;\ 0{,}1;\ 0{,}25)$, $\pd = 0{,}3$ ;
  \item salon qui se salit vite : même batterie, $\pd = 0{,}9$ ;
  \item très bonne batterie : $e_b = (0;\ 0{,}05;\ 0{,}5)$, $\pd = 0{,}3$.
\end{itemize}

Dans les trois cas, la politique optimale (identique par itération de la valeur et de la politique) est :
\begin{center}
\begin{tabular}{l c c c}
\toprule
& batterie haute & batterie moyenne & batterie basse \\
\midrule
Action & \N & \C & \C \\
\bottomrule
\end{tabular}
\end{center}
quels que soient la position du robot et l'état du salon.

\subsection{Interprétation}

\begin{itemize}
  \item \textbf{Le niveau de batterie devient le seul critère de décision.} Le risque de panne dépend désormais
        uniquement de la batterie, si bien que les états L$\ast b$ et C$\ast b$ ont exactement la même valeur.
        La position du robot n'apporte plus d'information.
  \item \textbf{Le robot nettoie uniquement batterie pleine.} En batterie moyenne, nettoyer coûte à la fois un
        risque de panne plus élevé et une descente en batterie basse, d'où il faudra de toute façon remonter.
        Charger ne coûte qu'un pas et ramène en batterie haute, où nettoyer est presque sans risque.
  \item \textbf{La politique est robuste.} Elle ne change ni avec $\pd$, ni avec une batterie haute parfaite :
        l'écart de risque entre les niveaux suffit à imposer le cycle « nettoyer à pleine charge, recharger dès que la
        batterie baisse ».
\end{itemize}

Ce modèle corrige la principale limite du modèle de base, où nettoyer depuis la station était toujours sans risque.

%==========================================================================
\section{Limites}
%==========================================================================

\begin{itemize}
  \item La charge est très avantageuse : un seul pas suffit pour gagner un niveau, sans aucun risque. Une charge
        probabiliste ou plus lente rendrait les allers-retours vers la station plus coûteux.
  \item Le salon ne se salit que pendant la charge ; en réalité, il se salit aussi pendant que le robot est dans le salon.
  \item Les politiques dépendent fortement du rapport entre $\mathit{costcrash}$, $r_c$ et $\gamma$. Une étude de
        sensibilité sur ces paramètres compléterait l'analyse.
\end{itemize}

\end{document}
